\chapter{Operações Lógicas}

As operações lógicas são fundamentais no estudo da lógica proposicional. Elas permitem a construção de proposições mais complexas a partir de proposições simples, obedecendo às regras do cálculo proposicional.

As principais operações lógicas são:
\begin{itemize}
    \item Negação;
    \item Conjunção;
    \item Disjunção;
    \item Disjunção exclusiva;
    \item Condicional;
    \item Bicondicional.
\end{itemize}

\section{Negação}

\subsection{Definição}
Chama-se \textbf{negação de uma proposição} $p$ a proposição representada por ``não $p$''.

\subsection{Notação}
A negação de $p$ é representada por:
\[
\sim p
\]

\subsection{Semântica}
A tabela-verdade da negação é dada por:

\begin{center}
\begin{tabular}{|c|c|}
\hline
$p$ & $\sim p$ \\
\hline
V & F \\
F & V \\
\hline
\end{tabular}
\end{center}

Ou seja, a negação inverte o valor lógico da proposição.

\subsection{Exemplos}
\begin{itemize}
    \item $p$: Santa Cruz é melhor do que o Sport.\\
    $\sim p$: Santa Cruz \textbf{não} é melhor do que o Sport.

    \item $q$: O professor de Lógica é o mais bonito da UFJ.\\
    $\sim q$: \textbf{É falso} que o professor de Lógica é o mais bonito da UFJ.
\end{itemize}

\subsection{Observação}
Se uma proposição é falsa, sua negação será verdadeira. Por exemplo:
\begin{itemize}
    \item $p$: 12 é um número ímpar (F)\\
    \item $\sim p$: Não é verdade que 12 é um número ímpar (V)
\end{itemize}

\section{Conjunção}

\subsection{Definição}
Chama-se \textbf{conjunção de duas proposições} $p$ e $q$ a proposição representada por ``$p$ e $q$''.

\subsection{Notação}
A conjunção é representada por:
\[
p \wedge q
\]

\subsection{Semântica}
A tabela-verdade da conjunção é:

\begin{center}
\begin{tabular}{|cc|c|}
\hline
$p$ & $q$ & $p \wedge q$ \\
\hline
V & V & V \\
V & F & F \\
F & V & F \\
F & F & F \\
\hline
\end{tabular}
\end{center}

A conjunção só é verdadeira quando ambas as proposições são verdadeiras.

\subsection{Exemplo}
\begin{itemize}
    \item $p$: Rivaldo foi revelado no Santa Cruz. \hspace*{0.5cm} (V)
    \item $q$: Neymar é melhor do que Cristiano Ronaldo. \hspace*{0.5cm} (F)

    \item $p \wedge q$: Rivaldo foi revelado no Santa Cruz e Neymar é melhor do que Cristiano Ronaldo.

    \item $V(p \wedge q) = V(p) \wedge V(q) = V \wedge F = F$
\end{itemize}

\subsection{Interpretação}
Para que uma conjunção seja verdadeira, todas as partes devem ser verdadeiras. Basta que uma seja falsa para que toda a proposição seja falsa.

\section{Disjunção}

\subsection{Definição}
Chama-se \textbf{disjunção de duas proposições} $p$ e $q$ a proposição representada por ``$p$ ou $q$''.

\subsection{Notação}
A disjunção é representada por:
\[
p \vee q
\]

\subsection{Semântica}
A tabela-verdade da disjunção é:

\begin{center}
\begin{tabular}{|cc|c|}
\hline
$p$ & $q$ & $p \vee q$ \\
\hline
V & V & V \\
V & F & V \\
F & V & V \\
F & F & F \\
\hline
\end{tabular}
\end{center}

A disjunção é falsa apenas quando ambas as proposições são falsas.

\subsection{Exemplo}
\begin{itemize}
    \item $p$: Rivaldo foi revelado no Santa Cruz. \hspace*{0.5cm} (V)
    \item $q$: Neymar é melhor do que Cristiano Ronaldo. \hspace*{0.5cm} (F)

    \item $p \vee q$: Rivaldo foi revelado no Santa Cruz ou Neymar é melhor do que Cristiano Ronaldo.

    \item $V(p \vee q) = V(p) \vee V(q) = V \vee F = V$
\end{itemize}

\subsection{Interpretação}
Na lógica proposicional, o ``ou'' é inclusivo, ou seja, a proposição é verdadeira se pelo menos uma das partes for verdadeira.

\section{Observações Finais}

As operações apresentadas neste capítulo — negação, conjunção e disjunção — são a base para a construção de expressões lógicas mais complexas.

Outras operações importantes, que serão estudadas posteriormente, incluem:
\begin{itemize}
    \item Disjunção exclusiva;
    \item Condicional;
    \item Bicondicional.
\end{itemize}

O domínio dessas operações é essencial para a compreensão de argumentos lógicos, provas matemáticas e aplicações em Ciência da Computação.